\documentclass[11pt]{article}
\usepackage[utf8]{inputenc}
\usepackage[T1]{fontenc}
\usepackage{geometry}
\usepackage{enumitem}
\usepackage{hyperref}
\usepackage{xcolor}
\usepackage{titlesec}
\usepackage{parskip}

\geometry{a4paper, margin=0.75in}
\hypersetup{colorlinks=true, linkcolor=blue, urlcolor=blue}
\titleformat{\section}{\large\bfseries\color{blue!70!black}}{\thesection}{1em}{}[\titlerule]

\begin{document}

\begin{center}
    {\LARGE\bfseries Multi-Agent GenAI Platform}\\[0.3em]
    {\large Cloud-Native AI System with RAG, LangGraph, and FastAPI}\\[0.5em]
    \href{https://your-live-url.a.run.app}{\texttt{Live Demo}} \textbar{} 
    \href{https://github.com/yourusername/multi-agent-genai}{GitHub Repository}
\end{center}

\section{Project Overview}
Production-style multi-agent AI system designed as a portfolio project to demonstrate cloud-native deployment, agent orchestration, and Retrieval-Augmented Generation (RAG). The platform exposes a FastAPI gateway serving both REST endpoints and an interactive web UI.

\section{Key Features}
\begin{itemize}[leftmargin=*]
    \item \textbf{FastAPI Gateway}: RESTful API with OpenAPI docs, rate limiting, and Bearer token auth
    \item \textbf{LangGraph Orchestration}: Multi-agent workflow with Planner, Retrieval, Tool Executor, Synthesizer, Critic, and Repair nodes
    \item \textbf{RAG Pipeline}: Document ingestion, chunking, embedding, and cosine-similarity search
    \item \textbf{Tool Execution Framework}: Document search, SQLite query, Python execution, and web search
    \item \textbf{Critic Agent}: Output validation with one automatic retry loop
    \item \textbf{Observability}: Structured JSON logging, in-memory metrics, and request tracing
\end{itemize}

\section{Technical Stack}
\begin{itemize}[leftmargin=*]
    \item \textbf{Backend}: Python 3.13, FastAPI, Uvicorn, Pydantic
    \item \textbf{AI/ML}: LangGraph, Groq API (Llama 3.1), OpenAI-compatible LLM client
    \item \textbf{Data}: SQLite, JSON vector store (MVP), pypdf for PDF parsing
    \item \textbf{Cloud}: Google Cloud Run, Container Registry, IAM Service Accounts
    \item \textbf{DevOps}: Docker, Cloud Build, GitHub Actions-ready CI/CD
    \item \textbf{Frontend}: Vanilla HTML/JS (single-file, no build step)
\end{itemize}

\section{Cloud Deployment}
\begin{itemize}[leftmargin=*]
    \item Deployed on \textbf{Google Cloud Run} with containerized FastAPI app
    \item Achieved \textbf{\$0 monthly infrastructure cost} by leveraging Cloud Run free tier (2M requests/month)
    \item Implemented multi-provider LLM client supporting Groq, Gemini, OpenAI, and local Ollama
    \item Externalized all secrets and configuration to \texttt{.env} and GCP environment variables
    \item Managed IAM permissions via service account keys with least-privilege roles
\end{itemize}

\section{Architecture Highlights}
\begin{itemize}[leftmargin=*]
    \item \textbf{StateGraph Workflow}: Planner $\rightarrow$ Retrieval $\rightarrow$ Tools $\rightarrow$ Synthesizer $\rightarrow$ Critic $\rightarrow$ [Retry $|$ End]
    \item \textbf{Cloud-Native}: Serverless container auto-scaling from 0 to N instances
    \item \textbf{Security}: Sandboxed Python executor, read-only SQL, optional Bearer auth
    \item \textbf{Extensibility}: Plugin-based tool registry and pluggable LLM providers
\end{itemize}

\section{Key Metrics}
\begin{itemize}[leftmargin=*]
    \item Multi-agent pipeline with 6 coordinated nodes
    \item 4 tool types (document, SQL, Python, web)
    \item Containerized deployment on GCP
    \item Live public URL for demo
\end{itemize}

\section{What I Learned}
\begin{itemize}[leftmargin=*]
    \item Designing agent workflows with LangGraph and state management
    \item Building production-ready FastAPI services with rate limiting and observability
    \item Deploying containerized Python apps to Google Cloud Run
    \item Managing cloud IAM roles, service accounts, and API enablement
    \item Securing API keys and credentials using environment variables
\end{itemize}

\end{document}
